\subsection{FOURIER TRANSFORMS OF INTEGRABLE FUNCTIONS}

In this number, we recall some definitions and results from Spectral Theories \(^{6}\) .

Denote by \(\hat{\mathbf{G}}\) the set of classes of irreducible representations of \(\mathbf{G}\) (on finite dimensional complex vector spaces). For all \(u\in \hat{\mathbf{G}}\) , denote by \(\mathrm{E}_u\) the space of \(u\) and \(d(u)\) its dimension. There exist separating positive hermitian forms on \(\mathrm{E}_u\) invariant under \(u\) , and any two such forms are proportional. Denote by \(\mathrm{A}^*\) (resp. \(\| \mathrm{A}\|_{\infty}\) ) the adjoint (resp. the norm) of an element \(\mathrm{A}\) of \(\operatorname {End}(\mathrm{E}_u)\) relative to one of these forms; for all \(g\in \mathbf{G}\) , we have \(u(g)^{*} = u(g)^{-1} = u(g^{-1})\) and \(\| u(g)\|_{\infty} = 1\) ; for all \(x\in \mathfrak{g}\) , we have \(u(x)^{*} = -u(x) = u(-x)\) .

Give \(\operatorname{End}(\mathbf{E}_u)\) the Hilbert space structure for which the scalar product is

\[
\langle \mathrm{A} | \mathrm{B} \rangle = d (u) \mathrm{Tr} (\mathrm{A} ^ {*} \mathrm{B}) = d (u) \mathrm{Tr} (\mathrm{BA} ^ {*}),\tag{1}
\]

and put

\[
\| \mathrm{A} \| _ {2} = \langle \mathrm{A} | \mathrm{A} \rangle^ {1 / 2} = (d (u) \mathrm{Tr} (\mathrm{A} ^ {*} \mathrm{A})) ^ {1 / 2}.\tag{2}
\]

We have

\[
\sqrt {d (u)} \| \mathrm{A} \| _ {\infty} \leq \| \mathrm{A} \| _ {2} \leq d (u) \| \mathrm{A} \| _ {\infty},\tag{3}
\]

SO

\[
| \langle \mathrm{A} | \mathrm{B} \rangle | \leq d (u) ^ {2} \| \mathrm{A} \| _ {\infty} \| \mathrm{B} \| _ {\infty}.\tag{4}
\]

For all \(g \in \mathrm{G}\) , we have \(\| u(g) \|_2 = d(u)\) .

Denote by \(\mathrm{F}(\hat{\mathrm{G}})\) the algebra \(\prod_{u\in \hat{\mathrm{G}}}\mathrm{End}(\mathrm{E}_u)\) . Denote by \(\mathrm{L}^2 (\hat{\mathrm{G}})\) the Hilbert sum of the Hilbert spaces \(\mathrm{End}(\mathrm{E}_u)\) ; this is the space of families \(\mathrm{A} = (\mathrm{A}_u)\in \mathrm{F}(\hat{\mathrm{G}})\) such that \(\sum_{u}\| \mathrm{A}_{u}\|_{2}^{2} < \infty\) , with the scalar product

\[
\langle \mathrm{A} | \mathrm{B} \rangle = \sum_ {u \in \hat {\mathrm{G}}} \langle \mathrm{A} _ {u} | \mathrm{B} _ {u} \rangle = \sum_ {u \in \hat {\mathrm{G}}} d (u) \operatorname{Tr} (\mathrm{A} _ {u} ^ {*} \mathrm{B} _ {u}).\tag{5}
\]

Denote the Hilbert norm on \(L^{2}(\hat{G})\) also by \(\parallel\parallel_{2}\) , so that \(\|A\|_{2}^{2}=\sum_{u\in\hat{G}}\|A_{u}\|_{2}^{2}\) for \(A\in L^{2}(\hat{G})\) .

If \(f\) is an integrable complex function on \(G\) , put

\[
u (f) = \int_ {\mathrm{G}} f (g) u (g) d g \in \operatorname{End} (\mathrm{E} _ {u})\tag{6}
\]

for all \(u \in \hat{\mathbf{G}}\) . We have \(\| u(f) \|_{\infty} \leq \int_{\mathbf{G}} |f(g)| dg = \| f \|_1\) . The Fourier cotransform of \(f\) , denoted by \(\overline{\mathcal{F}}(f)\) , is the family \((u(f))_{u \in \hat{\mathbf{G}}} \in \mathrm{F}(\hat{\mathbf{G}})\) . If \(f \in \mathrm{L}^2(\hat{\mathbf{G}})\) ,

\[
\| f \| _ {2} ^ {2} = \sum_ {u \in \hat {\mathsf {G}}} \langle u (f) | u (f) \rangle = \| \overline {{\mathcal {F}}} (f) \| _ {2} ^ {2},
\]

so \(\overline{F}\) induces an isometric linear map from the Hilbert space \(\mathrm{L}^{2}(\mathrm{G})\) to the Hilbert space \(\mathrm{L}^{2}(\hat{\mathrm{G}})\) : in other words, for f and \(f'\) in \(\mathrm{L}^{2}(\mathrm{G})\) , we have

\[
\int_ {\mathrm{G}} \overline {{f (g)}} f ^ {\prime} (g) d g = \langle \overline {{\mathcal {F}}} (f) | \overline {{\mathcal {F}}} (f ^ {\prime}) \rangle = \sum_ {u \in \hat {\mathrm{G}}} d (u) \mathrm{Tr} (u (f) ^ {*} u (f ^ {\prime})).\tag{7}
\]

For \(f\) and \(f'\) in \(\mathrm{L}^1(\mathrm{G})\) , the convolution product \(f * f'\) of \(f\) and \(f'\) is defined by

\[
(f * f ^ {\prime}) (h) = \int_ {G} f (h g ^ {- 1}) f ^ {\prime} (g) d g = \int_ {G} f (g) f ^ {\prime} (g ^ {- 1} h) d g
\]

(the integral makes sense for almost all \(h \in \mathbf{G}\) ).

We have \(f * f' \in L^{1}(G)\) and, for all \(u \in \hat{G}\) , \(u(f * f') = u(f)u(f')\) , so

\[
\overline {{{{\mathcal {F}}}}} (f * f ^ {\prime}) = \overline {{{{\mathcal {F}}}}} (f). \overline {{{{\mathcal {F}}}}} (f ^ {\prime}).\tag{8}
\]

Conversely, let \(\mathrm{A} = (\mathrm{A}_{u})_{u \in \hat{\mathrm{G}}}\) be an element of \(\mathrm{F}(\hat{\mathrm{G}})\) ; for all \(u \in \hat{G}\) , let \(F_{u}A\) be the (analytic) function on G defined by

\[
(\mathcal {F} _ {u} \mathrm{A}) (g) = \langle u (g) | \mathrm{A} _ {u} \rangle = d (u) \mathrm{Tr} (\mathrm{A} _ {u} u (g) ^ {- 1}).\tag{9}
\]

If \(\mathrm{A} \in \mathrm{L}^2(\hat{\mathrm{G}})\) , the family \((\mathcal{F}_u\mathrm{A})_{u \in \hat{\mathrm{G}}}\) is summable in \(\mathrm{L}^2(\mathrm{G})\) ; the Fourier transform of \(\mathrm{A}\) , denoted by \(\mathcal{F}(\mathrm{A})\) , is the sum of this family. The maps \(\overline{\mathcal{F}}\) and \(\mathcal{F}\) are inverse isomorphisms between the Hilbert spaces \(\mathrm{L}^2(\mathrm{G})\) and \(\mathrm{L}^2(\hat{\mathrm{G}})\) .

In other words:

PROPOSITION 1. Every square integrable complex function \(f\) on \(G\) is the sum in the Hilbert space \(L^2(G)\) of the family \((f_u)_{u\in \hat{G}}\) where, for all \(h\in G\) and all \(u\in \hat{G}\) ,

\[
\begin{array}{l} f _ {u} (h) = \langle u (h) | u (f) \rangle \\ \qquad = d (u) \int_ {\mathrm{G}} f (g) \mathrm{Tr} (u (g h ^ {- 1}))   d g = d (u) \int_ {\mathrm{G}} f (g h) \mathrm{Tr} (u (g))   d g. \end{array}\tag{10}
\]

For all \(u \in \hat{\mathrm{G}}\) choose an orthonormal basis \(\mathrm{B}_u\) of \(\mathrm{E}_u\) , and denote by \((u_{ij}(g))\) the matrix of \(u(g)\) in this basis. Prop. 1 also means that the family of functions \(\sqrt{d(u)} u_{ij}\) , for \(u\) in \(\hat{\mathrm{G}}\) and \(i,j\) in \(\mathrm{B}_u\) , is an orthonormal basis of the space \(\mathrm{L}^2 (\mathrm{G})\) .

If f is an integrable function on G such that the family \((f_{u})\) is uniformly summable, then the sum of this family is a continuous function which coincides almost everywhere with f; in other words, if we assume in addition that f is continuous, then for all \(h \in G\) ,

\[
f (h) = \sum_ {u \in \hat {\mathrm{G}}} d (u) \int_ {\mathrm{G}} f (g h) \mathrm{Tr} (u (g))   d g.\tag{11}
\]

Conversely, let \(\mathrm{A} \in \mathrm{F}(\hat{\mathrm{G}})\) ; if the family \((\mathcal{F}_u\mathrm{A})_{u \in \hat{\mathrm{G}}}\) is uniformly summable, the function

\[
g \mapsto \sum_ {u \in \hat {\mathrm{G}}} (\mathcal {F} _ {u} \mathrm{A}) (g) = \sum_ {u \in \hat {\mathrm{G}}} d (u) \mathrm{Tr} (\mathrm{A} _ {u} u (g) ^ {- 1})
\]

is a continuous function on \(G\) whose Fourier cotransform is \(A\) .

Let \(f\) be an integrable function on \(\mathbf{G}\) , and let \(s \in \mathbf{G}\) . Denote by \(\gamma(s)f\) and \(\delta(s)f\) the functions on \(\mathbf{G}\) defined by \(\gamma(s)f = \varepsilon_s * f\) , \(\delta(s)f = f * \varepsilon_{s^{-1}}\) , that is,

\[
(\gamma (s) f) (g) = f (s ^ {- 1} g), \quad (\delta (s) f) (g) = f (g s) \text { for } g \in \mathrm{G},
\]

(Chap. III, §3, no. 4 and Integration, Chap. VII, §1, no. 1). We have

\[
u (\gamma (s) f) = \int_ {\mathrm{G}} f (s ^ {- 1} g) u (g) d g = \int_ {\mathrm{G}} f (g) u (s g) d g,
\]

SO

\[
u (\gamma (s) f) = u (s) u (f),\tag{12}
\]

and similarly

\[
u (\delta (s ^ {- 1}) f) = u (f) u (s).\tag{13}
\]

When G is commutative, \(\hat{G}\) is the underlying set of the dual group of G (Spectral Theories, Chap. II, §1, no. 1), \(d(u) = 1\) for all \(u \in \hat{G}\) , and we recover the definitions of the Fourier transform given in Spectral Theories, Chap. II.

